\section{Acknowledgments}\label{sec:ack}

I thank Peter Kagey for posing the problem, and for discussing its early stages with me,
including the uniqueness of the $5\times5$ optimum and the first graph-theoretic approach to
the narrow grids. The design of Figures~\ref{fig:game} and~\ref{fig:boards} follows the figure
in his posting of Problem 001~\cite{kagey001}.

\section*{A note on the use of AI}

I began this work in January of 2026 and wrote the first draft and its mathematics
between August and October, together with most of the Python code. I used
Claude Opus 5.5 (Anthropic) for the October 17--19 revision, including the
cut of a draft from 45 to 30 pages and the language passes.
It did not add any language, just proposed a cut version.
Using this as reference, I rewrote the text myself and implemented the revisions as I saw fit.
No sentence in this paper was written by an LLM.
It also drew the figures in the style of Kagey's posting of the problem given my own hand-drawn sketches.

I used it to check my Python code, and it wrote further checks of the
computer-assisted results. One of these, \texttt{six\_row\_collar\_6x9\_check.py},
gives the lower bound in the second open question of Section~\ref{sec:open}.
It also wrote the Lean files, the LRAT certificate pipeline and the formal
checks described in Theorem~\ref{thm:A} and Section~\ref{sec:verify}.
The formalization did not change a statement of the paper or find an error
or missing hypothesis. It changed how several results are established,
as Section~\ref{sec:verify} describes.

I also used Claude Opus 5.5 to search the literature beyond the sources I
had found. Several references came from that search. I read every cited
source, including the parts supporting each citation. Later reviews with
Claude Opus 5.5 and OpenAI Codex on October 23rd pointed out problems that I checked and
corrected, such as typos and inconsistencies or redundancies.
I checked these revisions and implemented them all myself, and take full responsibility for the
contents of the paper.
